\chapter{Anforderungen und Plattformauswahl}
\label{chap:auswahl}

Dieses Kapitel überführt die in Kapitel~\ref{chap:stand-der-technik} identifizierte Forschungslücke
und die in Kapitel~\ref{chap:poc} gewonnenen praktischen Erkenntnisse in eine konkrete
Plattformwahl. Dazu werden zuerst die Anforderungen an eine Forschungs-Cyber-Range abgeleitet und zu
einem Kriterienkatalog verdichtet. Anschließend werden die Kandidaten anhand dieser Kriterien
verglichen und es wird eine begründete Auswahl getroffen.

\section{Anforderungen und Bewertungskriterien}
\label{sec:anforderungen-kriterien}

Aus Forschungsperspektive ergeben sich an eine Cyber Range Anforderungen, die über klassische
Trainingsfunktionen hinausgehen.

Zentrale Voraussetzung ist die Reproduzierbarkeit. Identische Szenarien müssen sich unter
kontrollierten Randbedingungen wiederholt ausführen lassen, damit Ergebnisse zwischen Werkzeugen,
Konfigurationen oder Methoden vergleichbar werden. Das betrifft sowohl die technische Bereitstellung
der Umgebung als auch die Nachvollziehbarkeit des Szenarioablaufs und der Messpunkte.

Eng damit verbunden ist ein hoher Automatisierungsgrad. Szenarien und Umgebungen sollten möglichst
einfach beschrieben und automatisiert ausgerollt werden können, etwa über Infrastructure-as-Code.
Ebenso wichtig ist die Erweiterbarkeit, also die Möglichkeit, neue Szenarien, zusätzliche Messpunkte
oder weitere Integrationen zu ergänzen. Für die Erhebung und Auswertung von Messdaten kommen in
modernen Umgebungen häufig Observability-Stacks zum Einsatz, etwa Prometheus und Grafana für
Metriken oder der Elastic Stack und Grafana Loki für Logs.

Ein weiteres Kernkriterium ist die Integrationsfähigkeit externer Werkzeuge. Da in dieser Arbeit
konkrete Cyber-Resilienz-Use-Cases umgesetzt und bewertet werden, muss die Plattform Schnittstellen
für Monitoring, Detection und Auswertung unterstützen. Resilienzbezogene Untersuchungen setzen zudem
voraus, dass sich das Systemverhalten über die Zeit beobachten lässt, was eine geeignete
Beobachtbarkeit über Logs, Metriken und Events erfordert.

Schließlich sind Betrieb und Governance zu berücksichtigen. Eine Forschungsumgebung muss isoliert
und sicher betrieben werden können und sollte Mehrbenutzerbetrieb mit Rollen- und Zugriffsmodellen
unterstützen. Praktisch relevant sind außerdem die Dokumentationsqualität, die Aktivität der
Community, das Lizenzmodell und der Bereitstellungsaufwand, da sie die Machbarkeit der Umsetzung
stark beeinflussen.

Aus diesen Anforderungen leitet sich der folgende Kriterienkatalog ab. Er bildet die
Bewertungsgrundlage für den späteren Vergleich und ist so gewählt, dass er die Ziele der Arbeit
abbildet, ohne eine bestimmte Plattform vorauszusetzen.

\textbf{K1 -- Reproduzierbarkeit und Wiederholbarkeit:}
Kann die Plattform identische Umgebungen und Szenarien in gleicher Konfiguration wiederholt
bereitstellen, etwa durch versionierbare Definitionen und deterministische Provisionierung?

\textbf{K2 -- Automatisierung und Lifecycle-Steuerung:}
Unterstützt die Plattform den Lifecycle von Umgebungen (Provisionierung, Start/Stop, Reset, Cleanup)
weitgehend automatisiert und mit geringem manuellem Betriebsaufwand?

\textbf{K3 -- Integrationsfähigkeit externer Werkzeuge:}
Erlaubt die Plattform die Einbindung zusätzlicher Tools innerhalb der Sandbox oder am Rand der
Umgebung (z.\,B.\ Sensoren, Agents, SIEM/Logging, Monitoring), sodass Vergleiche unter gleichen
Bedingungen möglich sind?

\textbf{K4 -- Beobachtbarkeit und Datenzugang:}
Stellt die Plattform Mechanismen bereit, um Mess- und Auswertungsdaten systematisch zu erfassen und
zugänglich zu machen (z.\,B.\ Logs, Metriken, Events, Netzwerkdaten) und diese für Analysen zu
exportieren oder weiterzuverarbeiten?

\textbf{K5 -- Zugriffsmöglichkeiten und Nutzerinteraktion:}
Bietet die Plattform praktikable Zugriffsmöglichkeiten auf die bereitgestellten Systeme (z.\,B.\ SSH
sowie optional GUI-/Web-Konsole/RDP/VNC) und unterstützt sie eine nutzerfreundliche Interaktion
während der Durchführung?

\textbf{K6 -- Benutzer-, Rollen- und Rechteverwaltung:}
Unterstützt die Plattform Mehrbenutzerbetrieb mit Rollen- und Gruppenkonzepten, nachvollziehbarer
Zugriffskontrolle und klaren Mechanismen für Authentisierung und Autorisierung?

\textbf{K7 -- Wartbarkeit, Dokumentation und Community:}
Ist die Plattform realistisch betreib- und wartbar (z.\,B.\ nachvollziehbarer Installations- und
Update-Prozess), gut dokumentiert und durch aktive Weiterentwicklung sowie Austausch- und
Support-Kanäle unterstützt?

\section{Vergleichende Bewertung}
\label{sec:vergleich}

Auf Basis der Kriterien K1--K7 und der praktischen Erprobung aus Kapitel~\ref{chap:poc} werden die
quelloffenen Kandidaten CyberRangeCZ, OCR und KYPO gegenübergestellt. Tabelle~\ref{tab:plattformvergleich-matrix}
fasst die Bewertung zusammen.

\begin{table}[H]
\centering
\renewcommand{\arraystretch}{1.2}
\begin{tabular}{p{5.2cm}ccc}
\hline
\textbf{Kriterium} & \textbf{CyberRangeCZ} & \textbf{OCR} & \textbf{KYPO} \\
\hline
\textbf{K1 Reproduzierbarkeit \& Wiederholbarkeit}      & \cmark & \cmark & \cmark \\
\textbf{K2 Automatisierung \& Lifecycle-Steuerung}       & \cmark & $\sim$ & \cmark \\
\textbf{K3 Integrationsfähigkeit externer Werkzeuge}     & $\sim$ & ?      & ?      \\
\textbf{K4 Beobachtbarkeit \& Datenzugang}               & \cmark & ?      & \cmark \\
\textbf{K5 Zugriffsmöglichkeiten \& Nutzerinteraktion}   & \cmark & ?      & \cmark \\
\textbf{K6 Benutzer-, Rollen- \& Rechteverwaltung}       & \cmark & \cmark & \cmark \\
\textbf{K7 Wartbarkeit, Dokumentation \& Community}      & \cmark & $\sim$ & \xmark \\
\hline
\end{tabular}
\caption{Vergleichsmatrix der Plattformen anhand der Kriterien K1--K7 (\cmark\ erfüllt, $\sim$ teilweise erfüllt, \xmark\ nicht erfüllt, ? offen).}
\label{tab:plattformvergleich-matrix}
\end{table}

Die Vergleichsmatrix in Tabelle~\ref{tab:plattformvergleich-matrix} zeigt, dass CyberRangeCZ in
mehreren für diese Arbeit zentralen Kriterien nachweisbar stark abschneidet. Insbesondere die
praktische Erprobung im PoC belegt einen hohen Grad an Automatisierung und Lifecycle-Steuerung (K2)
sowie praktikable Zugriffsmöglichkeiten und Nutzerinteraktion (K5). Auch der Mehrbenutzerbetrieb und
die Zugriffskontrolle (K6) sind durch die vorhandenen Plattformkomponenten und das Portal-Konzept
klar adressiert. Zusammen mit der umfangreichen Dokumentation und aktiven Weiterentwicklung (K7)
reduziert dies das Umsetzungsrisiko.

KYPO erfüllt konzeptionell ebenfalls wesentliche Anforderungen, insbesondere hinsichtlich
reproduzierbarer Umgebungen und automatisierter Bereitstellung (K1, K2). Es basiert jedoch auf einer
älteren technischen Architektur und wird im Projektkontext zunehmend durch CyberRangeCZ abgelöst.
Dadurch entsteht im Vergleich ein Nachhaltigkeitsrisiko (K7), das für eine längerfristige
Forschungsumgebung relevant ist. OCR verfolgt einen modularen Ansatz, jedoch
konnten zentrale Punkte in der Erprobung nicht in gleicher Tiefe praktisch verifiziert werden.
Entsprechend werden einzelne Kriterien als teilweise erfüllt oder offen bewertet. Das betrifft vor
allem die konkrete technische Bereitstellung einer vollständigen Übungskette im PoC sowie die daraus
ableitbare Beurteilung einzelner Betriebs- und Zugriffsaspekte. Zusätzlich wirkt die verfügbare
Dokumentation im Vergleich weniger umfassend und die Community-Aktivität ist weniger ausgeprägt, was
die Fehlersuche und den Wissenstransfer im Implementierungsprozess erschweren kann.

\section{Gewichtung und Auswahl der Plattform}
\label{sec:gewichtung-auswahl}

Die Auswahl der Plattform stützt sich auf den Kriterienkatalog aus
Abschnitt~\ref{sec:anforderungen-kriterien} und die vergleichende Bewertung in
Tabelle~\ref{tab:plattformvergleich-matrix}. Sie bildet die Grundlage für die Umsetzung der Use
Cases in den folgenden Kapiteln.

\subsection{Gewichtung der Kriterien}
\label{subsec:gewichtung}

Für das Forschungsvorhaben stehen Reproduzierbarkeit und Wiederholbarkeit (K1), Automatisierung und
Lifecycle-Steuerung (K2), Integrationsfähigkeit externer Werkzeuge (K3) sowie Beobachtbarkeit und
Datenzugang (K4) im Vordergrund, da sie die Vergleichbarkeit von Experimenten und die spätere
Evaluation von Cyber-Resilienz-Use-Cases direkt beeinflussen. Benutzer-, Rollen- und
Rechteverwaltung (K6) sowie Wartbarkeit, Dokumentation und Community (K7) sind ebenfalls wesentlich,
weil sie die langfristige Nutzbarkeit und den stabilen Betrieb einer Forschungsumgebung
beeinflussen. Zugriffsmöglichkeiten (K5) werden unterstützend berücksichtigt, da sie die praktische
Durchführbarkeit von Übungen und Experimenten maßgeblich erleichtern.

\subsection{Begründete Auswahl}
\label{subsec:auswahl}

Aus Tabelle~\ref{tab:plattformvergleich-matrix} ergibt sich CyberRangeCZ als die am robustesten
begründbare Option. Im Unterschied zu den Alternativen konnten zentrale Kriterien nicht nur anhand
von Dokumentation, sondern zusätzlich durch den erfolgreich umgesetzten PoC praktisch untermauert
werden. Damit liegen belastbare Hinweise vor, dass die Plattform den automatisierten Aufbau und
Betrieb von Umgebungen (K2) sowie einen praktikablen Zugriff und die Durchführung von
Trainingsinstanzen (K5) unterstützt. Darüber hinaus sprechen die aktuelle Weiterentwicklung, die
verfügbare Dokumentation und die Support- und Community-Strukturen (K7) sowie die vorhandenen
Mechanismen für Mehrbenutzerbetrieb und Zugriffskontrolle (K6) für ein geringeres Implementierungs-
und Betriebsrisiko.

KYPO wird als fachlich relevanter Vorgängeransatz eingeordnet, weist jedoch im Vergleich ein
erhöhtes Nachhaltigkeitsrisiko auf, da der Schwerpunkt der Weiterentwicklung auf CyberRangeCZ liegt.
OCR stellt zwar einen interessanten modularen Ansatz dar, konnte in der Erprobung jedoch nicht in
allen entscheidenden Punkten praktisch verifiziert werden. Entsprechend wird OCR nicht als primäre
Zielplattform gewählt, kann aber als Referenz sowie für eine mögliche spätere Erweiterung oder
Gegenüberstellung herangezogen werden.

Auf Basis der Kriteriengewichtung und der in der Erprobung gewonnenen Erkenntnisse wird daher
CyberRangeCZ als Plattform dieser Arbeit ausgewählt.
